\begin{proposition}[C.3.10 (Proved): reduction to a specified annular orbit map]
Let $\mathcal Q=X_{c_0}/\rho$ be the orbit space for the circle of
C.3.1. Choose a small cusp disc with coordinate $q$, and a smooth
radial function $\chi$ supported in that disc, equal to one on a
smaller disc. On its punctured annulus, let
$v_\gamma=\operatorname{pr}_{\delta}\Pi_{c_0}\widehat\gamma$
be the projected real period. Define an orbit map $\psi$ by
\[
 \psi([z])=[z+\chi(q)v_\gamma(q)]
\]
on the cusp annulus and by the identity elsewhere. This is a
well-defined diffeomorphism, equal to the identity near the fixed
set. It has a smooth equivariant lift to $X_{c_0}$. Every such lift
has mapping class $6[\Phi]$, for the positive $w$-loop convention
of C.1.1 and the closing map of C.3.5.

The assertion concerns the support of the \emph{orbit map}.
It does not assert an identity-ended lift supported in the annulus.
In particular it does not evaluate $6[\Phi]$.
\end{proposition}
\begin{proof}
We first specify the quotient transport used in the comparison.
Take $c(r)=c_0-12r$, $0\leq r\leq1$. On an ordinary fibre,
the quotient by $\rho$ is obtained by projecting the additive
period coordinates to $\mathbb C\times\mathbb R$, discarding
$\operatorname{Re}\zeta_2$. The period change
$(c(r)-c_0)e\otimes\gamma$ has real coefficient, so its
projected lattice is unchanged. This gives constant quotient
charts over the ordinary base.

The finite root quotient charts are also constant: the change of
their period columns and affine deck translations is in the real
$e$-direction and disappears under this projection. Their overlap
with the ordinary chart changes, however. With the notation and
signs of C.3.5,
\[
 \sigma_{j,c(r)}-\sigma_{j,c_0}
       =-12r\,\epsilon_j w_j e,
 \qquad w_j=(2\pi i)^{-1}\log s_j,
 \quad (\epsilon_1,\epsilon_2)=(1,-1).
\]
Let $R_a$ denote the positive real subgroup of the global
$\mathbb C^*$-action, with multiplier $e^a$. The transport in
the root chart which agrees with the ordinary identity is thus
$R_{12r\epsilon_j\log|s_j|}$ on the overlap. Choose a smooth
rotation-invariant real function $f_j(s_j)$ which vanishes near
zero and equals $\log|s_j|$ on that overlap. Replacing this by
$R_{12r\epsilon_j f_j(s_j)}$ extends the transport over the
root centre and commutes with its finite deck group. These
functions are smooth on the free root-cover quotient; no assertion
that $|t|^{1/m_j}$ is a smooth coordinate on the coarse base is
used. At the cusp, the varying multiplicative period unit is a
$\delta$-phase, so the toric quotient charts and their overlaps
are constant. We take the identity transport there.

This quotient transport is induced by a smooth equivariant
transport on the actual family. Here is the lifting justification.
On an ordinary or finite root chart the affine family and its
smooth lattice coordinates give local smooth horizontal lifts of
the prescribed quotient vector field. At the cusp use its smooth
toric charts, with the finite base parameter interval retained;
the dependence just removed is a circle phase, so these charts
give local equivariant lifts too, including at fixed points.
Two such lifts have the same quotient field and differ by a
vertical circle field on the free stratum. The difference vanishes
at a fixed point. An invariant smooth partition of unity on the
actual total space patches these local fields. Their parameter
component remains one and their quotient is the specified field.
Properness gives a flow for the whole compact interval. Denote
its endpoint by $H$. This argument uses smooth local lifts in the
given slices, rather than assuming an arbitrary orbit-space
isotopy lifts across the fixed set.

Put $G=\widetilde F_{12,c_0-12}\circ H$. By C.1.1 this smooth
equivariant diffeomorphism represents $6[\Phi]$. On ordinary
quotient charts, C.3.5 gives the positive real flow
$R_{a(t)}$, where
\[
 a(t)=\log|K(t)|,\qquad K(t)=t^4/(t-1)^3.
\]
On the finite root chart the same map is the flow with time
\[
 a_j(s_j)=\log|K(t)s_j^{-12\epsilon_j}|
                       +12\epsilon_j f_j(s_j).
\]
The first term is the logarithm of a nonvanishing unit and the
second is smooth; hence this time is smooth through the root
centre. On its outer overlap it equals $a(t)$. At the cusp the
quotient map is the bounded toric unit $U=U_{12,c_0-12}$ of
C.3.5. Write $u(q)$ for its additive projected logarithm.
That logarithm is single-valued on the full disc. The exact
period identity of C.3.5 is
\[
 u(q)+v_\gamma(q)=R_{a(q)}
\]
in the additive universal covering vector space; on the right
we mean the additive displacement of that real flow.

Multiplication by $U^{-s\chi(q)}$, $0\leq s\leq1$, defines a
smooth equivariant isotopy supported in the cusp. To see this
through every toric stratum, choose the explicit holomorphic
additive logarithm of the two nonzero unit coordinates of $U$.
Exponentiating $-s\chi$ times that logarithm gives a torus
element depending smoothly on $q$; multiplication acts smoothly
in every toric chart and commutes with the period identifications.
It is the identity near the outer edge of the disc.

There is a second global equivariant isotopy, the real
$\mathbb C^*$-flow with time $-s(1-\chi)a$, using $a_j$ in
the finite charts. The times agree on overlaps, are smooth on
the actual root quotients, and vanish on an inner cusp disc.
They are constant along the $\mathbb C^*$-orbits, so the inverse
flow is obtained by negating the time. Consequently this is an
isotopy of diffeomorphisms, not just a pointwise multiplication.

Compose $G$ with the endpoints of these two isotopies. Away
from the cusp its orbit map becomes the identity. In the cusp
annulus its projected displacement is exactly
\[
 a-\chi u-(1-\chi)a=\chi(a-u)=\chi v_\gamma.
\]
Where $\chi=1$, this is a whole projected lattice period and
the orbit map is again the identity; on the central fibre the
cusp-unit description gives the same identity directly. Thus the
result is an actual smooth equivariant lift of the stated $\psi$,
and has the same mapping class as $G$.

Finally $A_0\widehat\gamma=\widehat\gamma-\widehat\delta$
makes $v_\gamma$ single-valued on the punctured cusp after the
$\delta$-quotient. Translation by $\chi v_\gamma$ and its
negative are smooth mutually inverse maps on that free orbit
annulus, and they are the identity on collars of its ends.
This proves the direct well-definedness and smoothness assertion.
C.3.6 shows that every other existing smooth equivariant lift has
the same mapping class. This construction supplies no assertion
that the full lift is supported in the annulus: its residual
circle gauge outside that annulus has not been discarded.
\end{proof}
